\section*{序論}
8月19日朝までの24時間は、生成AIの性能競争そのものよりも、利用者層の拡大に伴う安全設計、frontier modelの訓練環境、著作権・透明性、そしてAI基盤のセキュリティが前面に出た。OpenAIでは13--17歳向けの専用体験と、最新モデル向け強化学習の一時停止が同じ観測窓に入り、製品展開と内部統制を並行して強める動きが注目された。

企業・制度面ではAnthropicの収益規模を巡る報道、Claude生成文への不可視透かし、ByteDanceと映画業界のIP保護枠組みが話題となった。一方、インフラ層ではRayの重大RCEがCISAの既知悪用脆弱性一覧に加わったとの警告が広がり、agentic commerceではAlipayがMCPを含む商取引向け基盤を打ち出した。全体として、生成AIを社会へ広げるための「能力の外側」の設計が主題となった一日である。
